% ==========================================
% Prerequisites (chapter 0), written by Claude (claude-opus-5-5) on 2026-10-11.
% Not part of the lectures: the background the lectures take for granted, with the notation of
% these notes. Results are stated as `fact` without proof, except Bézout's identity and Euclid's
% lemma, which get a short argument because the lectures build on them. The lectures point here
% with \cref{...}, which prints numbers 0.x.
% Extended on 2026-10-11 (same model): images and preimages, induction and the binomial theorem,
% the Euclidean algorithm, Hom-sets, and quotient vector spaces with the linear first isomorphism
% theorem, each because a lecture uses it without stating it.
% Extended on 2026-10-12 (same model), after a review of the lectures: the order of an element and
% cyclic groups, open sets and closure in R, smooth functions vanishing on a half-line, division
% with remainder for polynomials, and the notations N_0 and R^{N_0}.
% ==========================================

% Primary Language: English

\prerequisiteschapter{Notation and Facts Assumed by the Lectures}{Background from the First-Year Courses}

\begin{overview}[About This Chapter]
This chapter is not part of the lectures. It collects what the lectures take for granted from the first-year courses (Linear Algebra, Analysis, and the Foundations course \qt{Grundstrukturen}), together with the notation of these notes. The lectures point here wherever they use one of these facts.
\begin{itemize}
    \item Results are stated as \emph{facts}, without proof: they are standard first-year material. Bézout's identity and Euclid's lemma get a two-line argument, because the lectures build on them.
    \item The first lecture of the course (Tuesday, 15 September) was not recorded. According to the later lectures, it introduced groups, rings and fields and discussed factorization in $\mathbb{Z}$; these notes start with the second lecture. The lecturer restates groups and fields in Week~2 and rings in Week~1.
\end{itemize}
\end{overview}

\section{Notation, Sets and Maps}

\begin{notation}[Number Sets and Inclusions]\label[notation]{not:pre-notation}
\begin{itemize}
    \item $\mathbb{N} = \{0, 1, 2, \dots\}$, also written $\mathbb{N}_0$ to stress that $0$ is included, and $\mathbb{Z} \subset \mathbb{Q} \subset \mathbb{R} \subset \mathbb{C}$. Subscripts restrict a set: $\mathbb{Z}_{>0} = \{1, 2, 3, \dots\}$, $\mathbb{Z}_{\ge 2}$, $\mathbb{R}_{>0}$.
    \item For a field $K$, $K^* = K \setminus \{0\}$; for a ring, $R^*$ often denotes its units instead (Week~3, \cref{prop:units-mod-m-w3}).
    \item $A \subset B$ means that every element of $A$ lies in $B$; equality is allowed, as in the lecturer's usage. A proper inclusion is written $A \subsetneq B$.
    \item $|X|$ (or $\# X$) is the number of elements of a finite set $X$; for a finite group it is called the \newterm{order}.
    \item $X \times Y = \{ (x, y) \mid x \in X, \ y \in Y \}$ is the \newterm{Cartesian product}. A binary operation on $X$ is a map $X \times X \to X$; the set of pairs used for the field of fractions (\cref{def:set-of-pairs-s}) is a subset of $R \times R$.
    \item $Y^X$ is the set of all maps $X \to Y$. In particular $R^{\mathbb{N}_0}$ is the set of all sequences $(a_0, a_1, a_2, \dots)$ with entries in $R$; polynomials and power series are defined as such sequences (\cref{def:poly-rigorous-sequences}, \cref{def:power-series-ring-formal}).
\end{itemize}
\end{notation}

\begin{definition}[Maps]\label[definition]{def:pre-maps}
A map $f \colon X \to Y$ is \newterm{injective} if $f(x_1) = f(x_2)$ implies $x_1 = x_2$, \newterm{surjective} if every $y \in Y$ is of the form $f(x)$, and \newterm{bijective} if it is both. The composition of $f$ with $g \colon Y \to Z$ is $g \circ f \colon X \to Z$, $x \mapsto g(f(x))$; the identity of $X$ is $\id_X$ (also written $1_X$); the restriction of $f$ to $U \subset X$ is $f|_U \colon U \to Y$.
\end{definition}

\begin{fact}[Inverse Maps]\label[fact]{fact:pre-inverse-map}
A map $f \colon X \to Y$ is bijective if and only if there is a map $g \colon Y \to X$ with $g \circ f = \id_X$ and $f \circ g = \id_Y$. This $g$ is unique; it is the \newterm{inverse map} $f^{-1}$.
\end{fact}

\begin{definition}[Image and Preimage of a Subset]\label[definition]{def:pre-image-preimage}
Let $f \colon X \to Y$ be a map. For $A \subset X$, the \newterm{image} of $A$ is $f(A) = \{ f(a) \mid a \in A \} \subset Y$; the image of $f$ is $\operatorname{im}(f) = f(X)$. For $B \subset Y$, the \newterm{preimage} (or pullback) of $B$ is
\[
f^{-1}(B) = \{ x \in X \mid f(x) \in B \} \subset X.
\]
\end{definition}

\begin{importantremark}[Two Meanings of $f^{-1}$]\label[importantremark]{rem:pre-preimage-notation}
The preimage $f^{-1}(B)$ is defined for \emph{every} map $f$, bijective or not; it is a set, possibly empty. Only when $f$ is bijective does $f^{-1}$ also denote the inverse map (\cref{fact:pre-inverse-map}), and then the two readings agree: the preimage of $B$ is the image of $B$ under the inverse map. The kernel of a homomorphism is a preimage, $\ker(f) = f^{-1}(\{0\})$.
\end{importantremark}

\begin{fact}[Images and Preimages]\label[fact]{fact:pre-image-preimage}
Let $f \colon X \to Y$ be a map, $A, A' \subset X$ and $B, B' \subset Y$.
\begin{enumerate}[label=\textbf{(\roman*)}]
    \item Preimages respect inclusions, intersections and unions: $B \subset B' \implies f^{-1}(B) \subset f^{-1}(B')$, $f^{-1}(B \cap B') = f^{-1}(B) \cap f^{-1}(B')$, $f^{-1}(B \cup B') = f^{-1}(B) \cup f^{-1}(B')$. Images respect inclusions and unions, but in general not intersections.
    \item $f(f^{-1}(B)) \subset B$, with equality for every $B$ if and only if $f$ is surjective.
    \item $A \subset f^{-1}(f(A))$, with equality for every $A$ if and only if $f$ is injective.
\end{enumerate}
The Correspondence Theorem (\cref{prop:correspondence-theorem-ideals}, Week~4) uses (ii) for the surjective quotient map, and replaces injectivity in (iii) by the hypothesis that the ideal contains the kernel.
\end{fact}

\section{Equivalence Relations}

\begin{definition}[Equivalence Relation]\label[definition]{def:pre-equivalence-relation}
An \newterm{equivalence relation} on a set $T$ is a relation $\equiv$ on $T$ that is
\begin{enumerate}[label=\textbf{(\roman*)}]
    \item \textbf{reflexive:} $x \equiv x$ for all $x \in T$;
    \item \textbf{symmetric:} $x \equiv y \implies y \equiv x$;
    \item \textbf{transitive:} $x \equiv y$ and $y \equiv z \implies x \equiv z$.
\end{enumerate}
\end{definition}

\begin{definition}[Equivalence Classes and Quotient Set]\label[definition]{def:pre-quotient-set}
The \newterm{equivalence class} of $x \in T$ is $[x] = \{ y \in T \mid y \equiv x \}$, and the \newterm{quotient set} $T/{\equiv}$ (written $T_{\equiv}$ on the lecturer's board) is the set of all equivalence classes. Two classes are either equal or disjoint, $[x] = [y]$ holds exactly when $x \equiv y$, and the classes cover $T$: they form a \newterm{partition} of $T$.
\end{definition}

\begin{steps}[Maps Defined on Classes]
A rule that defines a value for a class through one of its elements (a \newterm{representative}) defines a map on $T/{\equiv}$ only if equivalent representatives give the same value; the rule is then called \newterm{well-defined}. The lectures carry out this check for the operations of $\mathbb{Z}/n\mathbb{Z}$ (\cref{def:residue-class-ring}), for the extension $g \colon \mathbb{Q} \to F$ (\cref{clm:g-well-defined}) and for the field of fractions (\cref{prop:addition-well-defined}).
\end{steps}

\section{The Integers}

\begin{fact}[Well-Ordering Principle]\label[fact]{fact:pre-well-ordering}
Every non-empty subset of $\mathbb{N}$ has a least element. (This is equivalent to the principle of induction, \cref{fact:pre-induction}.) The lectures use it to find the generator of an ideal of $\mathbb{Z}$ (\cref{thm:ideals-of-z-are-principal}).
\end{fact}

\begin{fact}[Induction]\label[fact]{fact:pre-induction}
Let $P(n)$ be a statement for each $n \in \mathbb{N}$. If $P(0)$ holds, and $P(n)$ implies $P(n+1)$ for every $n$, then $P(n)$ holds for all $n \in \mathbb{N}$. The same holds with any starting value $n_0 \in \mathbb{Z}$ in place of $0$.
\end{fact}

\begin{fact}[Binomial Theorem]\label[fact]{fact:pre-binomial}
For $n \in \mathbb{N}$ and $0 \le a \le n$, the \newterm{binomial coefficient} $\binom{n}{a} = \frac{n!}{a! \, (n-a)!}$ is a positive integer, and $\binom{n}{a} = \binom{n-1}{a-1} + \binom{n-1}{a}$ for $1 \le a \le n - 1$ (Pascal's rule). For all $x, y \in \mathbb{C}$,
\[
(x + y)^n = \sum_{a=0}^{n} \binom{n}{a} x^a y^{n-a},
\]
by induction on $n$ using Pascal's rule. The proof uses only commutativity and distributivity, so the formula holds in every commutative ring (\cref{thm:binomial-formula}, Week~4).
\end{fact}

\begin{fact}[Division with Remainder]\label[fact]{fact:pre-division}
For $y \in \mathbb{Z}$ and $n \in \mathbb{Z}_{>0}$ there are unique integers $q, r$ with
\[
y = q n + r \quad \text{and} \quad 0 \le r < n.
\]
\end{fact}

\begin{definition}[Divisibility, Greatest Common Divisor, Primes]\label[definition]{def:pre-divisibility}
For $a, b \in \mathbb{Z}$, $a$ \newterm{divides} $b$, written $a \mid b$, if $b = a c$ for some $c \in \mathbb{Z}$. For $a, b$ not both $0$, $\gcd(a, b)$ is the largest integer dividing both; the lecturer also writes $(a, b)$. The integers $a, b$ are \newterm{coprime} (relatively prime) if $\gcd(a, b) = 1$. A \newterm{prime} is an integer $p \ge 2$ whose only positive divisors are $1$ and $p$.
\end{definition}

\begin{fact}[Bézout's Identity]\label[fact]{fact:pre-bezout}
For $a, b \in \mathbb{Z}$, not both $0$, there are $x, y \in \mathbb{Z}$ with
\[
a x + b y = \gcd(a, b).
\]
In particular, $a$ and $b$ are coprime if and only if $a x + b y = 1$ for some $x, y \in \mathbb{Z}$.
\end{fact}

\begin{steps}[Why It Holds]
The set $I = \{ a x + b y \mid x, y \in \mathbb{Z} \}$ is an ideal of $\mathbb{Z}$ (\cref{def:ideal}), so $I = d\mathbb{Z}$ with $d > 0$ by \cref{thm:ideals-of-z-are-principal} (Week~2). Since $a, b \in I$, $d$ divides $a$ and $b$; and every common divisor of $a$ and $b$ divides every element of $I$, in particular $d$. Hence $d = \gcd(a, b)$. Conversely, if $a x + b y = 1$, every common divisor of $a$ and $b$ divides $1$.
\end{steps}

\begin{fact}[Euclid's Lemma]\label[fact]{fact:pre-euclid}
Let $p$ be a prime. If $p \mid a b$, then $p \mid a$ or $p \mid b$.
\end{fact}

\begin{steps}[Why It Holds]
If $p \nmid a$, then $\gcd(a, p) = 1$, since the only positive divisors of $p$ are $1$ and $p$. By Bézout, $a x + p y = 1$ for some $x, y \in \mathbb{Z}$, so $b = a b \, x + p \, b y$, and $p$ divides both terms on the right.
\end{steps}

\begin{example}[The Euclidean Algorithm]\label[example]{ex:pre-euclidean-algorithm}
The coefficients in Bézout's identity can be computed. Since $\gcd(y, n) = \gcd(n, r)$ whenever $y = q n + r$, repeated division with remainder (\cref{fact:pre-division}) shrinks the pair until the remainder is $0$; the last non-zero remainder is the gcd. For $a = 252$, $b = 198$:
\[
\begin{aligned}
252 &= 1 \cdot 198 + 54, &\qquad 198 &= 3 \cdot 54 + 36, \\
54 &= 1 \cdot 36 + 18, & 36 &= 2 \cdot 18 + 0,
\end{aligned}
\]
so $\gcd(252, 198) = 18$. Solving the equations backwards for the remainders,
\[
18 = 54 - 36 = 54 - (198 - 3 \cdot 54) = 4 \cdot 54 - 198 = 4 \cdot (252 - 198) - 198 = 4 \cdot 252 - 5 \cdot 198.
\]
If $\gcd(a, m) = 1$, the same back-substitution gives $a x + m y = 1$, so $\bar{x}$ is the inverse of $\bar{a}$ in $\mathbb{Z}/m\mathbb{Z}$ (\cref{thm:units-z-mod-m}).
\end{example}

\begin{fact}[Unique Factorization in $\mathbb{Z}$]\label[fact]{fact:pre-unique-factorization}
Every integer $n \ge 2$ is a product of primes, and this product is unique up to the order of the factors.
\end{fact}

\begin{definition}[Congruence Modulo $n$]\label[definition]{def:pre-congruence}
Let $n \in \mathbb{Z}_{>0}$. Integers $a, b$ are \newterm{congruent modulo $n$}, written $a \equiv b \pmod n$, if $n \mid a - b$.
\end{definition}

\begin{fact}[Residue Classes]\label[fact]{fact:pre-residue-classes}
Congruence modulo $n$ is an equivalence relation on $\mathbb{Z}$, and it is compatible with the operations: $a \equiv a'$ and $b \equiv b'$ imply $a + b \equiv a' + b'$ and $a b \equiv a' b' \pmod n$. Its classes are the \newterm{residue classes} $\bar{a} = a + n\mathbb{Z}$; by division with remainder there are exactly $n$ of them, $\bar{0}, \bar{1}, \dots, \overline{n-1}$. The lecturer recalls the resulting ring $\mathbb{Z}/n\mathbb{Z}$ in Week~1 (\cref{def:residue-class-ring}).
\end{fact}

\section{Groups, Fields and the Complex Numbers}

\begin{definition}[Group]\label[definition]{def:pre-group}
A \newterm{group} $(G, \cdot, e)$ is a set $G$ with an associative operation $\cdot$, a neutral element $e$ ($e a = a e = a$ for all $a$), and an inverse for every element ($a b = b a = e$ for some $b$). It is \newterm{abelian} if $a b = b a$ for all $a, b$; an abelian group is often written additively as $(G, +, 0)$, with inverses $-a$. Examples: $(\mathbb{Z}, +, 0)$ and $(\mathbb{R}^*, \cdot, 1)$. The lecturer restates the definition in Week~2 (\cref{def:group}).
\end{definition}

\begin{definition}[Group Homomorphism]\label[definition]{def:pre-group-hom}
A map $f \colon G \to H$ between groups is a \newterm{homomorphism} if $f(a b) = f(a) f(b)$ for all $a, b \in G$, and an \newterm{isomorphism} if it is moreover bijective (Week~2, \cref{def:group-homomorphism}).
\end{definition}

\begin{definition}[Order of an Element, Cyclic Groups]\label[definition]{def:pre-cyclic}
Let $G$ be a group with neutral element $e$ and $a \in G$. The \newterm{order} of $a$ is the least $n \ge 1$ with $a^n = e$ (or $\infty$ if there is none); in additive notation the condition reads $n a = 0$. The group $G$ is \newterm{cyclic}, with \newterm{generator} $a$, if every element of $G$ is a power $a^k$, $k \in \mathbb{Z}$. A finite group $G$ is cyclic with generator $a$ exactly when $a$ has order $|G|$, and then $G = \{e, a, a^2, \dots, a^{|G|-1}\}$. Example: $(\mathbb{Z}/n\mathbb{Z}, +, \bar{0})$ is cyclic with generator $\bar{1}$. An isomorphism maps elements of order $n$ to elements of order $n$; the lectures use this in Week~2 (\cref{prop:r-star-not-isom-r-pos}, \cref{prop:z7z-cyclic-isomorphism}).
\end{definition}

\begin{definition}[Field]\label[definition]{def:pre-field}
A \newterm{field} $(K, +, \cdot, 0, 1)$ is a set with two operations such that $(K, +, 0)$ is an abelian group, $(K \setminus \{0\}, \cdot, 1)$ is an abelian group, and $a (b + c) = a b + a c$ for all $a, b, c$. In particular $0 \neq 1$. Examples are $\mathbb{Q}$, $\mathbb{R}$, $\mathbb{C}$, and $\mathbb{Z}/p\mathbb{Z}$ for a prime $p$ (in the lectures: \cref{prop:z-mod-m-domain} and \cref{thm:units-z-mod-m}). The lecturer restates fields in Week~2 (\cref{def:field}).
\end{definition}

\begin{fact}[Complex Conjugation and Modulus]\label[fact]{fact:pre-complex}
For $z = a + b i \in \mathbb{C}$ ($a, b \in \mathbb{R}$, $i^2 = -1$), the \newterm{complex conjugate} is $\bar{z} = a - b i$ and the \newterm{modulus} is $|z| = \sqrt{a^2 + b^2}$. For all $z, w \in \mathbb{C}$:
\[
\overline{z + w} = \bar{z} + \bar{w}, \qquad \overline{z w} = \bar{z} \, \bar{w}, \qquad z \bar{z} = |z|^2, \qquad |z w| = |z| \, |w|,
\]
and $\bar{z} = z$ exactly when $z \in \mathbb{R}$. So conjugation respects sums and products, and under multiplication the lengths of complex numbers multiply (and their angles add).
\end{fact}

\section{Linear Algebra}

\begin{definition}[Linear Map, Kernel and Image]\label[definition]{def:pre-linear-map}
Let $V, W$ be vector spaces over a field $K$. A map $T \colon V \to W$ is \newterm{linear} if $T(u + v) = T(u) + T(v)$ and $T(\lambda v) = \lambda T(v)$ for all $u, v \in V$ and $\lambda \in K$. Its \newterm{kernel} is $\ker(T) = \{ v \in V \mid T(v) = 0 \}$ and its \newterm{image} is $\operatorname{im}(T) = T(V)$. The linear maps $V \to W$ form a vector space $\operatorname{Hom}_K(V, W)$ under pointwise operations; the lectures reuse the notation $\operatorname{Hom}$ for the morphisms of any category (\cref{def:category-foundations}).
\end{definition}

\begin{fact}[Kernel, Image and Rank--Nullity]\label[fact]{fact:pre-rank-nullity}
For a linear map $T \colon V \to W$, $\ker(T) \subset V$ and $\operatorname{im}(T) \subset W$ are subspaces, and $T$ is injective if and only if $\ker(T) = \{0\}$. If $V$ is finite-dimensional,
\[
\dim V = \dim \ker(T) + \dim \operatorname{im}(T).
\]
\end{fact}

\begin{fact}[Bases]\label[fact]{fact:pre-basis}
Every vector space spanned by finitely many vectors has a finite basis, and all its bases have the same number of elements, its \newterm{dimension}. That every vector space has a basis is proved with Zorn's lemma, which the course takes up itself (Week~1, \cref{q:vector-space-basis}).
\end{fact}

\begin{fact}[Division with Remainder for Polynomials]\label[fact]{fact:pre-poly-division}
Let $K$ be a field, $f, g \in K[x]$ polynomials with coefficients in $K$, and $g \neq 0$. There are unique $q, r \in K[x]$ with
\[
f = q g + r \quad \text{and} \quad r = 0 \text{ or } \deg(r) < \deg(g).
\]
This is long division of polynomials: as long as $\deg(f) \ge \deg(g)$, subtracting $\frac{a}{b} x^{\deg f - \deg g} \, g$, where $a$ and $b$ are the leading coefficients of $f$ and $g$, lowers the degree of $f$; dividing by $b$ is where $K$ has to be a field. It is the polynomial analogue of \cref{fact:pre-division}; the lectures use it for $g = x^2 + 1$ (Week~4, \cref{lem:poly-division-algorithm}).
\end{fact}

\begin{definition}[Quotient Vector Space]\label[definition]{def:pre-quotient-space}
Let $U \subset V$ be a subspace. The relation $v \equiv w \iff v - w \in U$ is an equivalence relation on $V$, with classes the cosets $v + U$. The \newterm{quotient space} $V/U$ is the set of these cosets with
\[
(v + U) + (w + U) = (v + w) + U, \qquad \lambda (v + U) = \lambda v + U;
\]
these operations are well-defined, $V/U$ is a vector space, and the projection $\pi \colon V \to V/U$, $v \mapsto v + U$, is linear and surjective with $\ker(\pi) = U$. If $V$ is finite-dimensional, $\dim V/U = \dim V - \dim U$.
\end{definition}

\begin{fact}[First Isomorphism Theorem for Linear Maps]\label[fact]{fact:pre-linear-first-iso}
For a linear map $T \colon V \to W$, the map $\bar{T} \colon V/\ker(T) \to \operatorname{im}(T)$, $v + \ker(T) \mapsto T(v)$, is a well-defined isomorphism of vector spaces. With $\dim V/\ker(T) = \dim V - \dim \ker(T)$ this gives rank--nullity again (\cref{fact:pre-rank-nullity}).
\end{fact}

\begin{steps}[The Model for Quotient Rings]
The quotient ring $R/I$ of Week~4 follows this pattern step by step, with the ideal $I$ in the role of the subspace $U$: congruence modulo $I$ (\cref{def:congruence-modulo-ideal}), the set of classes (\cref{def:quotient-set}), well-defined operations (\cref{prop:quotient-addition-well-defined}, \cref{prop:multiplication-well-defined}), and the First Isomorphism Theorem $R/\ker(f) \cong \operatorname{im}(f)$ (\cref{thm:first-isomorphism-theorem-statement}). The new ingredient is multiplication: a subspace only has to absorb scalars, an ideal has to absorb all of $R$, and that is exactly what makes the product of classes well-defined.
\end{steps}

\begin{definition}[Dual Space and Dual Map]\label[definition]{def:pre-dual-space}
The \newterm{dual space} of a vector space $V$ over $K$ is $V^* = \operatorname{Hom}_K(V, K)$, the vector space of linear maps $V \to K$. A linear map $T \colon V \to W$ has the \newterm{dual map} $T^* \colon W^* \to V^*$, $T^*(\varphi) = \varphi \circ T$; it satisfies $(S \circ T)^* = T^* \circ S^*$ and $(\id_V)^* = \id_{V^*}$.
\end{definition}

\begin{fact}[The Double Dual]\label[fact]{fact:pre-double-dual}
For every vector space $V$, the map $\operatorname{ev} \colon V \to V^{**}$, $v \mapsto (\varphi \mapsto \varphi(v))$, is linear and injective. If $V$ is finite-dimensional, then $\dim V^* = \dim V$ and $\operatorname{ev}$ is an isomorphism. (Used in Week~2, \cref{ex:double-dual-functor}.)
\end{fact}

\section{Analysis}

\begin{fact}[The Real Numbers]\label[fact]{fact:pre-reals}
$\mathbb{R}$ is a complete ordered field: every non-empty subset that is bounded above has a supremum, and every Cauchy sequence converges. It can be constructed from $\mathbb{Q}$ by Cauchy sequences or by Dedekind cuts (Week~3, Wednesday).
\end{fact}

\begin{fact}[Exponential Function and Logarithm]\label[fact]{fact:pre-exp}
The exponential function $\exp \colon \mathbb{R} \to \mathbb{R}_{>0}$ is a strictly increasing bijection with $\exp(0) = 1$ and $\exp(x + y) = \exp(x) \exp(y)$; its inverse is the natural logarithm $\ln \colon \mathbb{R}_{>0} \to \mathbb{R}$ (Week~2, \cref{prop:exp-isomorphism}).
\end{fact}

\begin{definition}[Open Sets and Closure in $\mathbb{R}$]\label[definition]{def:pre-open-closure}
A subset $U \subset \mathbb{R}$ is \newterm{open} if around each of its points it contains a whole interval $(x - \varepsilon, x + \varepsilon)$, $\varepsilon > 0$; examples are open intervals and $\mathbb{R}_{>0}$. A subset is \newterm{closed} if its complement is open. The \newterm{closure} $\overline{A}$ of $A \subset \mathbb{R}$ is the smallest closed set containing $A$: the set of all limits of convergent sequences in $A$. The lectures use these notions for the support of a function and for restricting functions to an open set (Week~2, \cref{def:compact-support-ideal}, \cref{def:restriction-hom}).
\end{definition}

\begin{fact}[Compact Sets and Smooth Functions]\label[fact]{fact:pre-compact}
A subset of $\mathbb{R}$ is \newterm{compact} if and only if it is closed and bounded (Heine--Borel); finite unions of compact sets, and closed subsets of a compact set, are compact. $C^\infty(\mathbb{R})$ denotes the infinitely differentiable functions $\mathbb{R} \to \mathbb{R}$; sums and products of such functions are again infinitely differentiable (Week~2, \cref{def:smooth-functions-ring}).
\end{fact}

\begin{fact}[Smooth Functions Vanishing on a Half-Line]\label[fact]{fact:pre-flat-function}
The function
\[
h(x) = \begin{cases} e^{-1/x} & x > 0, \\ 0 & x \le 0 \end{cases}
\]
lies in $C^\infty(\mathbb{R})$: all its derivatives tend to $0$ as $x \to 0^+$, because $e^{-1/x}$ goes to $0$ faster than any power of $x$. So $h$ vanishes exactly on $(-\infty, 0]$; the function $x \mapsto h(-x)$ vanishes exactly on $[0, \infty)$, and $h(x - a) \, h(b - x)$ is a smooth function that is non-zero exactly on $(a, b)$ (a \newterm{bump function}). Such functions show that the ideals of $C^\infty(\mathbb{R})$ in Week~2 are not zero and differ from each other (\cref{prop:compact-support-ideal}, \cref{prop:larger-ideal-smooth}).
\end{fact}

\begin{fact}[Geometric Series and Cauchy Product]\label[fact]{fact:pre-series}
For $z \in \mathbb{C}$ with $|z| < 1$,
\[
\sum_{k=0}^\infty z^k = \frac{1}{1 - z}.
\]
If $\sum_{n} a_n$ and $\sum_{n} b_n$ converge absolutely, then so does their \newterm{Cauchy product} $\sum_{n} c_n$ with $c_n = \sum_{i=0}^{n} a_i b_{n-i}$, and $\sum_n c_n = \bigl(\sum_n a_n\bigr)\bigl(\sum_n b_n\bigr)$. The lectures use the same coefficient formula to multiply polynomials and formal power series, where no convergence is involved (\cref{def:poly-ring-operations}, \cref{def:power-series-multiplication}).
\end{fact}
